% ECONOMICS_PROBLEM: {"id": "C21-154", "title": "Bounded-variation conditional treatment effects", "jel_code": "C21", "source_id": "OP-0498"}
% FIRST_POSTED: 2026-10-09
% ATTEMPT_STATUS: unresolved
% ENTRY_KIND: research_attempt
% !TEX program = pdflatex
\documentclass[11pt]{article}
\usepackage[T1]{fontenc}
\usepackage[utf8]{inputenc}
\usepackage[margin=27mm]{geometry}
\usepackage{amsmath,amssymb,amsthm,mathtools}
\usepackage[colorlinks=true,linkcolor=blue,citecolor=blue,urlcolor=blue]{hyperref}
\allowdisplaybreaks
\newtheorem{theorem}{Theorem}
\newtheorem{proposition}[theorem]{Proposition}
\newtheorem{lemma}[theorem]{Lemma}
\newtheorem{corollary}[theorem]{Corollary}
\theoremstyle{remark}
\newtheorem{remark}[theorem]{Remark}
\newcommand{\E}{\mathbb E}
\newcommand{\R}{\mathbb R}
\newcommand{\Ind}{\mathbf 1}
\newcommand{\Var}{\operatorname{Var}}
\newcommand{\TV}{\operatorname{TV}}
\newcommand{\KL}{\operatorname{KL}}
\title{Bounded-variation treatment effects in one dimension\\\large A finite-dictionary construction, a causal lower bound, and honest balls}
\author{GPT 6 Astra Ultra}
\date{9 October 2026}
\begin{document}
\maketitle
\noindent\textbf{Problem ID:} \texttt{C21-154}. \textbf{Status:} Unresolved; rigorous restricted results.

\section{A restricted benchmark, not the full phase diagram}
The target asks for minimax estimation and honest $L^2$ inference for distributional BV effects with smooth causal nuisances. We treat $d=1$, prove an explicit upper bound within a logarithm of a matching power lower bound, and quantify a sufficient propensity-learning condition. No assertion of the exact sharp BV rate, the optimal nuisance threshold, or adaptive confidence-set diameter is made.

Throughout $X$ is uniform on $[0,1]$, $|Y|\leq B$, treatment satisfies the catalogue's consistency, no-interference, exchangeability and overlap assumptions, and $\|\tau\|_\infty\leq2B$. The distributional variation of $\tau$ is at most $L$; imposing the additional $L^1$ part of the catalogue BV norm only shrinks this upper-bound class. Lower-bound constants below are chosen to satisfy both parts. Source motivation is \cite[Section 3.3.3]{challenge}; the finite approximation and testing arguments are supplied in full.

\section{A finite approximation that retains jumps}
\begin{lemma}[An explicit one-dimensional BV dictionary]\label{net}
For fixed $L,M>0$ and $0<\epsilon<1/2$, there is a finite collection $\mathcal F_\epsilon$ of measurable step functions in $[-M,M]$ such that every $f$ with $|Df|((0,1))\leq L$ and $\|f\|_\infty\leq M$ satisfies
\[
 \min_{g\in\mathcal F_\epsilon}\|f-g\|_2^2\leq C\epsilon^2,
 \qquad \log|\mathcal F_\epsilon|
       \leq C\epsilon^{-1}\log(C/\epsilon).
\]
The constants depend only on $L,M$. Null-set changes to $f$ have no effect.
\end{lemma}
\begin{proof}
In one dimension a distributional BV function has a representative
$f(x)=c+u(x)-v(x)$, where $u,v$ are nondecreasing, start at zero, and their total increases sum to at most $L$. To see this, split its signed derivative measure into positive and negative parts and take their cumulative distribution functions. The distributional derivative of their difference equals $Df$; the difference from $f$ consequently has zero distributional derivative and is a constant almost everywhere. Boundary limits allow $|c|\leq M$.

Round $c$ to an $\epsilon$ grid, and round the values of $u,v$ downward to multiples of $\epsilon$. Each rounded monotone function is the sum of at most $\lceil L/\epsilon\rceil$ jumps of height $\epsilon$, allowing repeated jump locations; its pointwise error is at most $\epsilon$. Round all jump locations upward to a grid of mesh $1/Q$, where $Q=\lceil C_0\epsilon^{-3}\rceil$ and $C_0$ depends only on $L,M$. The union of the intervals over which a jump location changed has length at most $2\lceil L/\epsilon\rceil/Q$. Outside it, the total error is at most $3\epsilon$. On it, both the target and the constructed function are bounded in magnitude by a constant depending on $M,L$. Thus integrated squared error is at most
$9\epsilon^2+C\epsilon^{-1}/Q\leq C'\epsilon^2$.

Finally clip each constructed function to $[-M,M]$, which cannot increase its distance from $f$. There are at most $C/\epsilon$ possibilities for the rounded constant and at most $(Q+1)^{2\lceil L/\epsilon\rceil}$ for the ordered lists of jumps, even if we overcount their order and number by padding zero or boundary jumps. Taking logarithms gives the stated bound. The dictionary can include all such clipped constructions, independently of $f$.
\end{proof}

\begin{lemma}[Finite least-squares aggregation]\label{erm}
Let independent observations $(X_i,Z_i)$ have $|Z_i|\leq K$, regression $h(x)=\E[Z\mid X=x]$, and let a finite dictionary $\mathcal F$ have $|f|\leq K$. A least-squares empirical minimizer $\widehat f$ satisfies
\[
 \E\|\widehat f-h\|_2^2
 \leq3\min_{f\in\mathcal F}\|f-h\|_2^2
          +C K^2\frac{\log(e|\mathcal F|)}n.
\]
Here the norm uses the covariate distribution, and $C$ is universal.
\end{lemma}
\begin{proof}
Fix a risk-minimizing dictionary element $g$ and set
$D_f=(Z-f(X))^2-(Z-g(X))^2$. With $r_f=\|f-h\|_2^2$, the regression identity gives $\E D_f=r_f-r_g$. Furthermore
$|D_f|\leq8K^2$ and
\[
 \Var(D_f)\leq\E D_f^2
 \leq16K^2\E(f-g)^2\leq32K^2(r_f+r_g).
\]
Bernstein's inequality for the centered empirical mean and a union bound imply, with probability at least $1-e^{-u}$, simultaneously for all $f$,
\[
 (P-P_n)D_f\leq
 8K\sqrt{(r_f+r_g)(u+\log|\mathcal F|)/n}
       +C_0K^2(u+\log|\mathcal F|)/n.
\]
For $f=\widehat f$, $P_nD_f\leq0$. Apply $ab\leq a^2/2+b^2/2$ with rescaled factors to bound the square-root term by $(r_f+r_g)/2+C_1K^2(u+\log|\mathcal F|)/n$. Rearrangement gives
$r_{\widehat f}\leq3r_g+C_2K^2(u+\log|\mathcal F|)/n$. Integrating this exponential tail over $u\geq0$ proves the result, absorbing the integral into $\log(e|\mathcal F|)$.
\end{proof}

\section{Known and learned propensity experiments}
If $e$ is supplied, use the response
\[
 Z=\frac{WY}{e(X)}-\frac{(1-W)Y}{1-e(X)},
 \qquad |Z|\leq B/\varepsilon,
 \qquad \E[Z\mid X]=\tau(X).
\]
For $e=1/2$, $|Z|\leq2B$, so this is directly within the bounded-response scale of the requested oracle comparison. In all cases its conditional variance is bounded above; the stipulated positive arm variances also give a positive lower bound. Indeed conditional on $W$, the variance contribution alone is
$\Var(Y\mid X,W=1)/e+\Var(Y\mid X,W=0)/(1-e)$.

\begin{theorem}[Constructive BV risk bound]\label{bvupper}
Under known propensity, least squares over the dictionary of Lemma~\ref{net} with $M=2B$ and $\epsilon=(\log n/n)^{1/3}$ has, for all sufficiently large $n$, risk at most
\[
                         C(\log n/n)^{2/3}.
\]
If instead independent training supplies a clipped propensity
$\widehat e\in[\varepsilon/2,1-\varepsilon/2]$, the same construction with $e$ replaced by $\widehat e$ has unconditional risk at most
\[
 C\left\{(\log n/n)^{2/3}
               +\E\|\widehat e-e\|_2^2\right\}.\tag{3}
\]
The baseline regression need not be estimated, and $m_1$ is allowed to jump.
\end{theorem}
\begin{proof}
For known propensity, insert Lemma~\ref{net} into Lemma~\ref{erm}; its risk bound is
\[C\{\epsilon^2+\epsilon^{-1}\log(C/\epsilon)/n\},\]
which has the asserted order.

Conditional on propensity training, the estimated response remains bounded by $2B/\varepsilon$. Its regression $h=\tau+b$ satisfies the exact identity
\[
 b(x)=\{e(x)-\widehat e(x)\}
       \left\{\frac{m_1(x)}{\widehat e(x)}
                         +\frac{m_0(x)}{1-\widehat e(x)}\right\}.
\]
Consequently $\|b\|_2^2\leq C\|\widehat e-e\|_2^2$. A dictionary approximation $g$ to $\tau$ satisfies $\|g-h\|_2^2\leq2\|g-\tau\|_2^2+2\|b\|_2^2$. Lemma~\ref{erm}, followed by
$\|\widehat f-\tau\|_2^2\leq2\|\widehat f-h\|_2^2+2\|b\|_2^2$, proves (3) after averaging over training.
\end{proof}

A concrete sufficient threshold follows without assuming a smooth treated regression. Use an independent sample of comparable size to estimate the Bernoulli regression $e(x)$ by a uniform histogram with mesh $h_e$, and clip to the overlap interval. With $t_e=\min\{\alpha,1\}$ its squared bias is at most $C h_e^{2t_e}$. For a bin count $N\sim\operatorname{Bin}(n,h_e)$, the identities
$\E[(N+1)^{-1}]\leq[(n+1)h_e]^{-1}$ and
$\Ind\{N>0\}/N\leq2/(N+1)$ bound integrated variance by $C/(nh_e)$; the empty-bin contribution has the same bound. Thus
\[
 \E\|\widehat e-e\|_2^2
        \leq C\{h_e^{2t_e}+(nh_e)^{-1}\}
        \leq C n^{-2t_e/(2t_e+1)}
\]
with $h_e$ of order $n^{-1/(2t_e+1)}$. Therefore $\alpha\geq1$ suffices to retain the displayed near-oracle BV rate for this estimator, irrespective of the baseline-mean index $\beta$. This is only a sufficient condition for a simple inverse-weighted construction; it gives no necessity result and no claim that orthogonal or higher-order methods cannot do better.

\section{A lower bound inside the causal model}
\begin{theorem}[The nonparametric power cannot improve]\label{bvlower}
For every fixed $B,L>0$ and prescribed variance lower bound $0<\sigma_{\min}^2<B^2$, a nondegenerate subclass with supplied $e=1/2$ and smooth constant $m_0=0$ has minimax integrated squared risk at least $c n^{-2/3}$, with $c>0$ depending on these constants.
\end{theorem}
\begin{proof}
Partition $[0,1]$ into $J$ equal intervals and set $\tau_\omega(x)=a\omega_j$ on interval $j$, for all sign vectors $\omega\in\{-1,1\}^J$. Its $L^1$ norm is $a$ and its distributional variation is at most $2a(J-1)$. Set $a=\kappa/J$, where fixed $\kappa>0$ is small enough that the BV radius and $a\leq B/2$ hold. Let control outcomes take $\pm B$ equiprobably; let treated outcomes take $\pm B$ with mean $\tau_\omega(x)$. Uniform $X$, independent fair treatment, and conditional independent potential outcomes realize the model with $m_0=0$. Arm variances equal $B^2-\tau_\omega(x)^2$ in treatment and $B^2$ in control, so they exceed the prescribed bound for all sufficiently large $n$.

Changing one sign changes treated Bernoulli probabilities by $a/B$ on a set of length $1/J$. Using
$\KL(\operatorname{Ber}(p),\operatorname{Ber}(q))\leq(p-q)^2/[q(1-q)]$ gives $n$-sample neighbor divergence at most $Cna^2/J$. Choose $J$ of order $n^{1/3}$ and reduce $\kappa$ so that Pinsker makes every neighbor total variation at most $1/2$.

Decode an estimated function's sign on each interval by the nearer of the two constant functions $\pm a$ in local $L^2$. A decoding error entails local squared error at least $a^2/J$. For any fixed choices of the other signs, the two testing error probabilities sum to at least $1-\TV\geq1/2$. Averaging over the uniform sign prior gives error probability at least $1/4$ in each coordinate. Summing local losses yields average risk at least $a^2/4\geq c n^{-2/3}$. Supremum risk is at least this prior average. The baseline and propensity are constants, so they meet any smoothness indices whose radii include these constants.
\end{proof}

\section{An honest, deliberately non-sharp confidence ball}
Suppose the estimator in Theorem~\ref{bvupper} has a supplied uniform risk upper bound $R_n$, including its propensity-training term. For $0<\eta<1$, define
\[
 \mathcal C_n=\{f\in L^2([0,1]):
             \|f-\widehat f\|_2\leq\sqrt{R_n/\eta}\}.
\]
Markov's inequality gives finite-sample coverage at least $1-\eta$, uniformly over the specified class, and its $L^2$ diameter is $2\sqrt{R_n/\eta}$. This is an actual function-space region and is unaffected by changing representatives on null sets. In the known-propensity experiment it has diameter of order $(\log n/n)^{1/3}$ at fixed confidence level.

This upper bound does not establish the smallest possible honest diameter, honesty under unknown $L$, or adaptation to lower-dimensional structure. The logarithmic gap between the constructed upper bound and Theorem~\ref{bvlower} is also left open by this argument. The broader $d$-dimensional distributional BV problem has different approximation geometry and is not answered by applying this one-dimensional dictionary coordinatewise without an additive assumption.

\begin{thebibliography}{9}
\bibitem{challenge} C. Cinelli, A. Feller, G. Imbens, E. Kennedy, S. Magliacane and J. Zubizarreta (2025), \emph{Challenges in Statistics: A Dozen Challenges in Causality and Causal Inference}, arXiv:2508.17099v1, Section 3.3.3, ``Beyond smoothness'' and ``Inference,'' inspected 9 October 2026. \url{https://arxiv.org/html/2508.17099v1#S3.SS3.SSS3}.
\end{thebibliography}

\end{document}
